% ---------------------------------------------------------------------
%  05-notation.tex -- the HPB notation standard (condensed)
%  Glyphs in tab:notation are drawn with hpb.sty v1.0 (same code as the
%  figures). The anatomy figure (fig:anatomy) is input in section 03.
% ---------------------------------------------------------------------
\section{The HPB Notation}
\label{sec:notation}

A methodology spreads when people can draw it on a whiteboard and others
read the drawing the same way. Git-flow~\citep{driessen2010gitflow} is a
good example: one diagram, copied everywhere, and every engineer knows how
to read it. HPB needs the same thing, so the notation is fixed: the canvas,
where each element goes, the glyphs and a short list of layout rules. All
diagrams in this paper are drawn with the \LaTeX{} package
\texttt{hpb.sty} (\cref{app:package}), and the open-source renderer of
\cref{sec:tooling} copies the same glyphs, so a diagram looks the same in
a paper, a wiki or a product portal. A quick-reference card is in
\cref{app:card}.

Four goals shaped the notation:
\begin{enumerate}[label=\textbf{G\arabic*},leftmargin=*,itemsep=1pt,topsep=3pt]
  \item \textbf{One concept, one glyph.} Every element a reader must
        recognize has exactly one glyph, and no glyph is decorative.
  \item \textbf{Position carries meaning.} \emph{When} and \emph{what kind
        of work} are read from horizontal and vertical position, before
        color or shape.
  \item \textbf{Never color alone.} Every category is encoded at least
        twice, so diagrams survive grayscale printing and color-vision
        deficiency.
  \item \textbf{Familiar to engineers.} A horizontal main line, version tags
        above the nodes, branches that leave and rejoin it: the vocabulary
        of the branch diagrams R\&D teams already read.
\end{enumerate}

\subsection{Canvas, lanes and reading order}
\label{sec:notation-canvas}

Time runs left to right. The vertical axis carries no quantity; it is
divided into horizontal \emph{lanes} of equal height. Lane~$0$ is the
\emph{main line}: a thick (3\,pt) dark stroke that \emph{is} the product.
Released baselines sit on it as nodes, in order; planned baselines sit to
the right of the present moment on a dashed continuation of the same
stroke. Each lane has a faint dotted guide across the canvas and a legend
box at its left end that names the type in words; the legend boxes and the
product box of the line align in one column, so the left margin reads as
the key.

\paragraph{Lane placement.}
Branches run on the lane of their type, by a fixed \emph{lane map}:
\begin{equation}
  \lambda(\mathsf{I}) = +1,\qquad
  \lambda(\mathsf{L}) = +2,\qquad
  \lambda(\mathsf{X}) = -1,\qquad
  \lambda(\mathsf{T}) = -2 .
  \label{eq:lanemap}
\end{equation}
Work that makes the product better (improvements, lab evolutions) goes
\emph{above} the main line; work that keeps it correct (issues, tasks)
goes \emph{below}. On each side, the type that most directly changes what
customers receive sits nearest the line. When several branches of a type
are open at the same time, extra lanes of that type are stacked outward
next to the existing ones, and the order $\mathsf{I}$ before $\mathsf{L}$
above and $\mathsf{X}$ before $\mathsf{T}$ below is kept
(\cref{fig:scale}). When a diagram shows several lines or streams, each
gets its own \emph{lane band}, separated by a thin rule, and above or below
is read relative to that band's main line
(\cref{fig:hardware-fork,fig:backport,fig:shared}).

\paragraph{Reading order.}
Every HPB diagram reads the same way: \emph{above the line, how the product
gets better; below it, what is wrong and how it is fixed; left of NOW,
history, which does not change; right of NOW, the plan, which is expected
to change.} The NOW marker is a dashed vertical line in its own color,
labeled \textsf{NOW}, across all lanes. Left of it all strokes are solid,
right of it all strokes are dashed: the drawing shows directly that
released versions are immutable. \Cref{fig:anatomy} annotates every glyph
of a single-line diagram.

\subsection{Ordinal time}
\label{sec:notation-time}

By default HPB uses an \emph{ordinal} time axis: every event (a release, a
branch start, a step, a terminal) gets its own column of equal width, in
time order, followed by the NOW column and then the planned baselines in
order of target date (\cref{def:ordinal}). Calendar months appear only as
tick labels on an axis below the diagram.

I chose ordinal time because real product histories are lumpy. A single
week can bring a hotfix, a release and three new issues, and then nothing
happens for a month. On a calendar scale the busy week becomes an
unreadable knot and the quiet month wastes half the page. With one column
per event, dense periods stay readable, glyphs never collide, and labels
can be placed by a fixed rule. The \emph{linear} alternative, where
horizontal position is proportional to calendar time as in a Gantt chart,
is still allowed when the distance to a date is itself the message, as in
a roadmap; it must then be stated in the caption.

\begin{remark}
  \label{rem:ordinal}
  In an ordinal diagram, horizontal length is \emph{not} duration. A long
  branch is one that was open while many other things happened, not
  necessarily one that was open for long. Durations such as lead time are
  computed from the model (\cref{sec:metrics}) and, when needed, printed
  as labels; they are never read off the drawing.
\end{remark}

\subsection{Glyphs}
\label{sec:notation-glyphs}

\Cref{tab:notation} lists the core glyphs and the model element each one
stands for; they are drawn with the same macros as the figures, so the
table is an exact sample of the notation. A few rules complete it.

\emph{Baselines.} When the target date of a planned baseline has passed, an
exclamation badge appears and its tag and date turn red. A hardware
baseline is a diamond with an indigo tag; with the caption \emph{new ref.}\
it is the first node of a new line, reached from the parent line by a
vertical lineage connector. A process baseline is a square. A maintenance
stream leaves its baseline with a rounded corner and continues in its own
band. Life-cycle milestones are thin labeled ticks, and a bar labeled
\textsf{EoS} ends a retired line. The dot at the left end of a main line
marks the creation of the reference; it is not a baseline.

\emph{Branches and terminals.} A branch is an orthogonal path with rounded
corners: a short run from the fork baseline, a drop (or rise) to the lane
of its type, and a run along the lane. Its label, above the start node,
gives the tracker reference and a short title. A discarded branch turns
gray, never rejoins the line and carries its rationale. A re-targeted
branch keeps its old target as a faded dashed continuation, so the change
of plan stays visible (\cref{fig:discard}).

\emph{Markers.} A step after closure, such as an effectiveness check on a
corrective action, sits on a dotted continuation (\cref{fig:rma}). A dark
badge under a baseline counts merged branches whose drawing was collapsed.
A dash-dotted frame around planned baselines of several lines is a launch
bundle. Thin straight connectors are the only lines that may cross lanes
diagonally: dotted with an open arrowhead for a dependency, dashed with a
filled arrowhead for a propagation.

% ---------------------------------------------------------------------
% Glyph table.  Glyph column: small hpbgraph pictures (hpb.sty v1.0).
% ---------------------------------------------------------------------
% \HPBglyph[<hpbgraph keys>]{<hpb macros>}: one glyph sample with 3pt of
% vertical padding (glyph sizes are normative; only spacing is reduced).
\providecommand{\HPBglyph}[2][]{%
  \begin{hpbgraph}[col=0.55cm,lane=0.9cm,legend=1.3cm,#1]#2%
    \path ([yshift=3pt]current bounding box.north)
          ([yshift=-3pt]current bounding box.south);%
  \end{hpbgraph}}

\begin{table}[p]
  \centering
  \small
  \caption{The HPB notation: glyph, element and meaning, drawn with
    \texttt{hpb.sty} at normative size. Color is never the only carrier of
    meaning (\cref{sec:notation-a11y}). The full card, with the macro that
    draws each glyph, is in \cref{app:card}.}
  \label{tab:notation}
  \renewcommand{\arraystretch}{1.12}
  \renewcommand{\tabularxcolumn}[1]{m{#1}}
  \begin{tabularx}{\linewidth}{@{}>{\centering\arraybackslash}m{3.55cm}
      >{\raggedright\arraybackslash}m{2.75cm}
      >{\raggedright\arraybackslash}X@{}}
    \toprule
    Glyph & Element & Meaning \\
    \midrule
    % main line
    \HPBglyph{\hpbmainline{0}{2.2}\hpbmainline*{2.2}{4.2}
      \draw[hpb/now] (2.2,0.35) -- (2.2,-0.35);}
    & Main line
    & The product line and its versions. Solid left of NOW (history),
      dashed right of NOW (plan). \\
    % origin
    \HPBglyph{\hpbmainline{0}{2}\hpborigin{0}{Product A}}
    & Line origin
    & Creation of the reference (not a version), with the product
      reference in a box aligned with the lane legends. \\
    % released baseline
    \HPBglyph{\hpbmainline{-1.2}{1.2}\hpbrelease{0}{v1.1}}
    & Released baseline
    & A released version; the tag holds its name. Immutable. \\
    % planned baseline
    \HPBglyph{\hpbmainline*{-1.8}{1.2}\hpbplanned{-0.7}{v1.2}{Mar 2027}}
    & Planned baseline
    & A planned version with its target date, and a gate. Badge ``!'',
      red tag and date: target date passed. \\
    % hardware baseline
    \HPBglyph{\hpbmainline{-1.8}{1.2}\hpbhardware{-0.7}{Rev B}}
    & Hardware baseline
    & A hardware revision; with the caption \emph{new ref.}, the first
      version of a new line (line fork). \\
    % gate band
    \HPBglyph[lane=0.75cm]{\hpbmainline*{1.3}{4.6}
      \node[hpb/planned] at (4,0) {};
      \hpbgate[0.8]{4}{0}{-1}{}
      \hpbterminal{target}{iss}{1.5}{-1}
      \hpbtarget{1.5}{-1}{4}}
    & Gate band
    & The targeted set $\mathcal{B}(g)$: every dashed continuation entering
      the band must be verified before the version ships
      (\cref{rule:gate}). \\
    % start nodes
    \HPBglyph[col=0.75cm]{\hpbstart{imp}{0}{0}{I}\hpbstart{lab}{1}{0}{L}
      \hpbstart{iss}{2}{0}{X}\hpbstart{task}{3}{0}{T}}
    & Start node
    & Opening of a branch; fill color and letter give its type. \\
    % branch path from a baseline
    \HPBglyph[lane=0.75cm]{\hpbmainline{-0.6}{3.8}
      \node[hpb/baseline] at (0,0) {};
      \hpbbranch{iss}{0}{-1}{1.3}{3.8}{none}{}}
    & Branch (fork)
    & The path leaves its fork baseline $f(b)$ and runs on the lane of its
      type. \\
    % merge
    \HPBglyph[lane=0.75cm]{\hpbmainline{-0.6}{4.6}
      \node[hpb/baseline] at (0,0) {};
      \node[hpb/baseline] at (4,0) {};
      \hpbbranch{imp}{0}{1}{1.2}{2.7}{merge}{4}}
    & Merge terminal
    & Merged: the path rejoins the version $m(b)$ that closed it. \\
    % discard
    \HPBglyph[lane=0.75cm]{\hpbmainline{-0.6}{3.8}
      \node[hpb/baseline] at (0,0) {};
      \hpbbranch{task}{0}{-1}{1.3}{3.3}{discard}{}}
    & Discard terminal
    & Discarded: closed without a version, with a mandatory rationale;
      gray path, never rejoins. \\
    % open head
    \HPBglyph[lane=0.75cm]{\hpbmainline{-0.6}{3.8}
      \node[hpb/baseline] at (0,0) {};
      \hpbbranch{lab}{0}{1}{1.3}{3.3}{open}{}
      \draw[hpb/now] (3.3,1.6) -- (3.3,0.4);}
    & Open branch
    & Open: the head lies on the NOW marker. \\
    % targeted
    \HPBglyph[lane=0.75cm]{\hpbmainline{-0.6}{2.7}\hpbmainline*{2.7}{5}
      \node[hpb/baseline] at (0,0) {};
      \node[hpb/planned] at (4.4,0) {};
      \hpbbranch{iss}{0}{-1}{1}{2.7}{target}{4.4}}
    & Targeted branch
    & Committed: dashed continuation from the head to the planned target
      $g(b)$. \\
    % steps
    \HPBglyph[col=0.6cm,lane=0.7cm]{\draw[hpb/path] (-0.4,0) -- (4.4,0);
      \hpbstep{imp}{0}{0}{done}\hpbstep{imp}{1.33}{0}{todo}
      \hpbstep{lab}{2.67}{0}{pass}\hpbstep{lab}{4}{0}{fail}}
    & Steps
    & Done, pending, test or verification passed, test failed. \\
    % NOW
    \HPBglyph[lane=0.7cm]{\hpbmainline{-1.6}{0}\hpbmainline*{0}{1.6}\hpbnow{0}{0}{0.2}}
    & NOW marker
    & The present moment: history to its left, plan to its right. \\
    \bottomrule
  \end{tabularx}
\end{table}

\subsection{Never color alone}
\label{sec:notation-a11y}

Color is lost in grayscale prints and photocopies, and many readers have
a color-vision deficiency. Color universal design therefore recommends
never using color as the only carrier of
information~\citep{okabe2008cud}. HPB encodes every category at least
twice:
\begin{itemize}[leftmargin=1.5em,itemsep=1pt,topsep=3pt]
  \item \emph{branch type}: lane position, letter in the start node, legend
        box in words, and fill color;
  \item \emph{branch state}: terminal shape (check, cross, hollow head),
        line style (solid, dashed, gray) and position relative to NOW;
  \item \emph{baseline status}: filled or white node, solid or dashed
        outline, presence of a target date, position relative to NOW;
  \item \emph{baseline kind}: shape (circle, diamond, square), tag color,
        and the caption \emph{new ref.}\ for a line fork;
  \item \emph{overdue}: the ``!'' badge in addition to the red tag;
  \item \emph{test outcome}: check or cross in addition to green or red.
\end{itemize}
The palette uses light fills with one dark ink color for every outline,
path and text. All fills keep a contrast ratio of at least $4.4{:}1$
against the ink, so letters, checks and crosses stay readable in
grayscale. Not every pair of fills is distinguishable without color
vision: task green and lab teal are close. That is acceptable only because
the two types are never adjacent (one is below the main line, the other
above it) and their letters and legend boxes differ. In the spirit of
Tufte~\citep{tufte2001}, there are no gridlines other than the faint lane
guides, no separate legend panel, and the thickest stroke is reserved for
the product itself.

\subsection{Layout rules}
\label{sec:notation-layout}

Glyphs alone do not make diagrams comparable; their arrangement must be
fixed too. A diagram \emph{conforms} to the HPB notation when it uses the
glyphs of \cref{tab:notation} and follows these rules. Every figure in
this paper conforms, \texttt{hpb.sty} implements the rules by construction
or by documented convention, and the renderer of \cref{sec:tooling}
applies them automatically, with the one gap noted there (the lane
preference of \ref{lay:onelane}).

\begin{enumerate}[label=\textbf{L\arabic*},ref=L\arabic*,leftmargin=*,
    itemsep=2pt,topsep=3pt]
  \item \label{lay:ordinal}\textbf{Ordinal columns.} Unless the caption
        says linear time: one event per column, equal column width, and one
        canvas (column width, lane height, key width) for all diagrams of a
        document. One exception is allowed: the merge terminals of the
        branches that a single release closes may share one column, just
        before that release, as long as they sit on different lanes. Calendar periods appear only as ticks below the diagram,
        and the axis may be left out when dates are not the message.
  \item \label{lay:lanes}\textbf{Fixed lanes.} Branches run on the lanes
        of \cref{eq:lanemap}; extra lanes of a type stack outward next to
        that type. Every lane has a guide and a legend box, aligned in the
        key column with the origin boxes of all lines.
  \item \label{lay:onelane}\textbf{One branch per lane at a time.} A lane
        is reused only after the previous branch on it has ended. Branches
        are placed in order of their fork. A branch already closed by NOW
        takes the innermost free lane of its type; a branch still open or
        targeted at NOW takes the outermost one.
  \item \label{lay:fork}\textbf{Clean verticals.} A branch leaves its fork
        baseline half a column to the right of the node and enters its
        merge baseline half a column to the left, so verticals never cross
        a tag. When several branches share a baseline and a side, the
        farther lane uses the smaller offset, so corners nest and paths
        never cross.
  \item \label{lay:tags}\textbf{One tag per baseline.} Directly above the
        node, with an arrow to it. Tags never overlap each other, a path or
        a label: make the column at least as wide as the widest tag, or
        insert a spacer column between consecutive baselines. In linear
        time, a colliding tag is raised one tag height.
  \item \label{lay:labels}\textbf{Labels.} A branch label (reference and
        short title) is left-aligned at the start node, above the lane, on
        a white halo, and ends before the next vertical and before NOW. The
        only text beside a terminal is the discard rationale. Nothing is
        written on a main line, and callouts never cover a path or a label
        other than the one they point to.
  \item \label{lay:terminal}\textbf{Exactly one terminal.} Every branch
        ends in exactly one merge, discard or open head, mirroring
        \cref{rule:closure}.
  \item \label{lay:now}\textbf{History is not redrawn.} Left of NOW,
        strokes are solid and drawn as recorded; right of NOW, only planned
        versions, targeted continuations and pending steps appear, all
        dashed or hollow. An audited reopen is a new event in its own
        column, never an edit of the past.
  \item \label{lay:linefork}\textbf{Several lines and streams.} Each child
        line, maintenance stream or line of a configuration gets its own
        lane band, separated by a thin rule, with its legend boxes in the
        key column. A child line starts at its hardware diamond, reached by
        the lineage connector; a stream leaves its fork baseline. A retired
        parent line ends with an \textsf{EoS} bar.
  \item \label{lay:fonts}\textbf{Legibility.} Scale diagrams uniformly,
        never in one direction only. After scaling, no text is smaller than
        6\,pt and glyphs keep their relative sizes.
\end{enumerate}

Rules \ref{lay:ordinal}--\ref{lay:onelane} also make the layout
\emph{computable}: columns come from sorting the events, lanes from
packing the branch intervals of each type. That is the basis of the
layout algorithm in \cref{app:layout}, and the reason a text description
of a product history is enough to produce a conforming diagram.
